\subsection{The Concrete Internal C Structure of the CPython Runtime}

CPython should not be imagined as a mysterious interpreter floating above the machine. It is a large native C program. It contains C source files, C structures, C functions, native memory allocation, operating-system calls, and a large amount of bookkeeping code whose purpose is to make Python's high-level behavior possible.

For this course, we do not need to read CPython as a core developer would read it. We do not need every source file, every macro, every optimization, or every C API function. But we do need the basic structural picture, because many Python behaviors become clear only when we understand what CPython is storing and moving underneath the surface.

The most important idea is this:

\begin{quote}
In CPython, Python-level concepts are represented by C-level structures.
\end{quote}

A Python object is represented by a C object structure. A Python type is represented by a C type object. A Python function has a runtime object that refers to executable code. A module has a namespace dictionary. A frame stores execution state for a running Python code block. Namespaces are dictionaries or dictionary-like mappings. The interpreter itself keeps runtime state, interpreter state, and thread state.

Thus, Python's apparent simplicity is built from a concrete internal object world.

\subsubsection{CPython Is Not One Object but a Runtime System}

There is no single C structure called ``the Python program.'' CPython is a system of cooperating structures. At a simplified level, it contains:

\begin{itemize}
    \item process-level runtime state;
    \item interpreter-level state;
    \item thread-level execution state;
    \item Python object structures;
    \item type objects;
    \item code objects;
    \item frame execution structures;
    \item namespace dictionaries;
    \item allocator and garbage-collection structures;
    \item import-system state;
    \item native extension interfaces.
\end{itemize}

These layers correspond to different questions. Runtime state answers: what does the CPython process know globally? Interpreter state answers: what belongs to this Python interpreter instance? Thread state answers: what is this executing thread currently doing? Object structures answer: what values exist? Frame structures answer: what code is currently executing? Namespace dictionaries answer: which names are bound to which objects?

This is why CPython cannot be reduced to ``an interpreter loop.'' The interpreter loop is central, but it is only the execution site. The loop needs objects to operate on, code objects to read from, frames to store execution state, dictionaries to resolve names, and type objects to decide how operations behave.

\subsubsection{The Object Header: Why Every Object Can Be Treated as an Object}

At the lowest practical level of Python values, CPython uses a common object header. Conceptually, an ordinary Python object begins with information like this:

\begin{verbatim}
PyObject
    reference count
    pointer to type object
\end{verbatim}

This is not meant as exact source-code replacement for all builds and all versions. It is the conceptual skeleton. The reference count helps CPython track object lifetime. The type pointer tells CPython what kind of object this is and which operations the object supports.

This is the concrete meaning behind the phrase ``everything is an object'' in CPython. It does not mean that objects are magical. It means that runtime values are represented as structured objects that CPython can access through a common initial layout.

For example, when CPython has a pointer to some runtime value, it can ask two fundamental questions:

\begin{itemize}
    \item How many strong references are currently keeping this object alive?
    \item Which type object describes this object's behavior?
\end{itemize}

The second question is especially important. CPython does not need the C compiler to know that a value is an integer, string, list, dictionary, function, or user-defined instance at every source location. The object itself carries a pointer to its runtime type.

\subsubsection{Variable-Sized Objects: The Role of \texttt{PyVarObject}}

Some Python objects have a natural size that varies from instance to instance. A tuple may have three elements or ten elements. A bytes object may contain five bytes or five thousand bytes. A string may contain different amounts of character data.

For such objects, CPython uses an extended base form commonly represented by \texttt{PyVarObject}. Conceptually:

\begin{verbatim}
PyVarObject
    PyObject base
    size field
\end{verbatim}

The important idea is simple. A variable-sized object still begins like an object, so CPython can treat it as a Python object. But it also stores size-related information needed by that particular family of objects.

This explains why Python containers are not C arrays in the naive sense. A Python list, tuple, string, or bytes object is not just raw adjacent user values. It is an object with a header, type information, size information, and implementation-specific storage.

\subsubsection{Type Objects: Behavior Stored as Runtime Metadata}

The type pointer inside each object points to a type object. In CPython, a type object is itself a large C-level structure. It stores information about the type's name, size, construction behavior, destruction behavior, attribute access, numeric operations, sequence operations, mapping operations, method tables, inheritance relationships, and other behavior.

A simplified view is:

\begin{verbatim}
object
    ob_type  --->  type object
                       name
                       size information
                       deallocation function
                       attribute access functions
                       call behavior
                       numeric operation slots
                       sequence operation slots
                       mapping operation slots
                       method definitions
\end{verbatim}

This is the implementation-level reason operators and built-in functions work dynamically. When Python evaluates an operation such as

\begin{verbatim}
a + b
\end{verbatim}

CPython does not treat \texttt{+} as one fixed machine instruction. It must look at the objects involved and use their type information to decide what addition means. Integer addition, string concatenation, list concatenation, NumPy array addition, and user-defined \texttt{\_\_add\_\_} behavior are not the same operation internally.

The type object is therefore the behavior descriptor of Python objects. It is what turns a memory object into a value with operations.

\subsubsection{Reference Counts: Object Lifetime as Runtime Bookkeeping}

CPython primarily tracks object lifetime through reference counting. When a strong reference to an object is created, the object's reference count is increased. When a strong reference is released, the count is decreased. When the last strong reference is released, CPython can call the object's deallocation function. The official C API exposes this through operations such as \texttt{Py\_INCREF()} and \texttt{Py\_DECREF()}, with the latter invoking the type's deallocation function when the last strong reference is gone. :contentReference[oaicite:1]{index=1}

This gives CPython a very direct memory-management style compared with tracing-only garbage collectors. Many objects can be destroyed as soon as they become unreachable through ordinary reference paths.

However, the student should not oversimplify this into ``Python immediately frees everything.'' There are complications:

\begin{itemize}
    \item some objects may be cached or reused;
    \item some objects may be immortal in modern CPython builds;
    \item reference cycles need garbage-collector support;
    \item finalizers can execute Python code during destruction;
    \item alternative Python implementations may not use the same lifetime strategy.
\end{itemize}

The useful mental model is not that memory disappears magically. The useful model is that CPython constantly maintains object-lifetime bookkeeping while Python code runs.

\subsubsection{Namespaces as Dictionaries of Object References}

Python names are not C variables holding raw values. Names live in namespaces. In CPython, the most important namespaces are represented by dictionaries or dictionary-like internal structures.

A module namespace is essentially a mapping from names to objects:

\begin{verbatim}
module namespace
    "x"       ---> integer object
    "message" ---> string object
    "func"    ---> function object
    "MyClass" ---> class object
\end{verbatim}

This is why assignment in Python is binding, not typed storage mutation. When Python executes

\begin{verbatim}
x = 10
\end{verbatim}

CPython creates or reuses an integer object and binds the name \texttt{x} to that object in the relevant namespace.

When Python later executes

\begin{verbatim}
x = "hello"
\end{verbatim}

the namespace entry for \texttt{x} is changed to point to a different object. There was no fixed \texttt{int} box that became a \texttt{str} box. There was a name-to-object binding that changed.

This one idea removes many beginner misconceptions about Python variables.

\subsubsection{Code Objects and Function Objects}

Python code is also represented by objects. When CPython compiles source code, it produces code objects. A code object is not a running function call. It is an executable description: bytecode, constants, names, local-variable information, and other metadata needed to execute a block of Python code.

A function object wraps a code object together with additional runtime context, especially the global namespace and default arguments.

A simplified view is:

\begin{verbatim}
function object
    code object
        bytecode
        constants
        names
        local variable metadata
    globals dictionary
    defaults
    closure cells, if needed
\end{verbatim}

This distinction matters. A function definition does not merely mark a region of source text. Executing a \texttt{def} statement creates a function object. That function object can be assigned, passed as an argument, returned from another function, stored in a list, or attached to a class.

This is the implementation basis for Python's first-class functions.

\subsubsection{Frame State: Running Code Needs a Runtime Record}

A code object is executable, but it is not by itself an active execution. To run code, CPython needs execution state. This state is represented by frame machinery.

A Python frame contains the information needed to execute a code object at a particular moment. Conceptually, it includes:

\begin{itemize}
    \item the code object being executed;
    \item the current instruction position;
    \item local variables or local reference slots;
    \item access to global and built-in namespaces;
    \item an evaluation stack for intermediate expression values;
    \item exception-handling state and control-flow metadata.
\end{itemize}

This is why tracebacks can show a chain of Python calls. CPython has structured runtime records for Python-level execution. When an exception propagates, CPython can report which frames were active and which lines or instructions were involved.

Again, this should not be confused with the native C stack alone. CPython itself uses the native C stack because it is a C program. But Python-level execution also has its own frame model. That frame model is what makes Python-level introspection, debugging, tracebacks, generators, and coroutines possible.

\subsubsection{Interpreter State, Thread State, and the GIL}

Above individual objects and frames, CPython also maintains larger execution state. The official C API documentation separates topics such as interpreter initialization, thread states, the global interpreter lock, and multiple interpreters in one process. :contentReference[oaicite:2]{index=2}

For our purposes, the simplified model is enough:

\begin{verbatim}
CPython process
    runtime state
        interpreter state
            modules
            builtins
            import machinery
            memory-management state
            thread states
                current frame / execution state
                exception state
                thread-specific runtime information
\end{verbatim}

The thread state connects a native thread to the Python interpreter machinery. CPython needs to know which thread is currently executing Python code, what frame it is executing, and what exception state belongs to it.

The Global Interpreter Lock, usually called the GIL, is part of this execution-control story in standard CPython builds. Its practical consequence is that CPython historically allows only one thread at a time to execute Python bytecode in one interpreter, even though threads may still be useful for I/O-bound work or for native code that releases the lock. The detailed history of the GIL belongs elsewhere, but its existence is important because it shows that Python execution is controlled by runtime state, not only by operating-system threads.

\subsubsection{Native Extensions Use the Same Object World}

One reason CPython became so important is that native extension modules can participate in this object system. A C extension can create Python objects, receive Python objects, inspect types, call functions, raise exceptions, define new types, and expose C or C++ functionality through Python interfaces.

This is how much of the high-performance Python ecosystem works. NumPy, parts of pandas, image-processing libraries, cryptographic libraries, database drivers, and many other packages cross the boundary between Python and native code.

The bridge works because CPython has a concrete C-level object protocol. A native extension does not receive vague Python values. It receives pointers to CPython object structures and obeys CPython's reference-counting and type rules.

This is powerful, but it also explains why many packages are CPython-specific. If a package depends on the CPython C API, then it is not merely depending on Python syntax. It is depending on CPython's concrete runtime representation.

\subsubsection{The Practical Mental Model}

The practical mental model is this:

\begin{verbatim}
Python source code
    becomes code objects

code objects
    execute inside frames

frames
    use namespaces to resolve names

namespaces
    map names to objects

objects
    have headers, reference counts, and type pointers

type objects
    define behavior through runtime metadata and C function slots

CPython
    coordinates all of this inside a native process
\end{verbatim}

This is enough internal structure for an effective Python programmer. The goal is not to memorize CPython's source tree. The goal is to understand why Python behaves differently from C.

Python assignment is name binding because namespaces map names to object references. Python operations are dynamic because objects carry type pointers. Python memory management is automatic at the source level because CPython maintains reference counts and garbage-collection support. Python function calls are introspectable because execution uses frame state. Python extension libraries can be fast because they cross into native code while still speaking CPython's object protocol.

The phrase ``Python is high-level'' therefore has a concrete meaning. It does not mean there is no machinery. It means the machinery has been moved into CPython.